\section{UC Berkeley Integration Bee}
\subsection*{2023}
\subsubsection*{Qualifying Exam}

\begin{questions}

\question
\[ \int_1^2 \frac{\{x\}^2}{x} \, dx \]
\begin{solution}
Recall that $\{x\} = x - \lfloor x \rfloor$ denotes the fractional part of $x$. For $x \in [1, 2)$, we have $\lfloor x \rfloor = 1$, so $\{x\} = x - 1$.
Substitute this into the integral:
\[ \int_1^2 \frac{(x-1)^2}{x} \, dx = \int_1^2 \frac{x^2 - 2x + 1}{x} \, dx = \int_1^2 \left( x - 2 + \frac{1}{x} \right) \, dx \]
Evaluate the antiderivative:
\[ \left[ \frac{x^2}{2} - 2x + \ln|x| \right]_1^2 = (2 - 4 + \ln 2) - \left( \frac{1}{2} - 2 \right) = \ln 2 - \frac{1}{2} \]
\end{solution}

\question
\[ \int \frac{1}{\cosh(\log(x)) + x} \, dx \]
\begin{solution}
Express $\cosh(\log(x))$ using exponential definitions:
\[ \cosh(\log(x)) = \frac{e^{\log(x)} + e^{-\log(x)}}{2} = \frac{x + \frac{1}{x}}{2} = \frac{x^2 + 1}{2x} \]
The denominator becomes:
\[ \cosh(\log(x)) + x = \frac{x^2 + 1}{2x} + x = \frac{3x^2 + 1}{2x} \]
Thus, the integral simplifies to:
\[ \int \frac{2x}{3x^2 + 1} \, dx \]
Let $u = 3x^2 + 1$, then $du = 6x \, dx \implies 2x \, dx = \frac{1}{3} du$:
\[ \frac{1}{3} \int \frac{1}{u} \, du = \frac{1}{3} \ln|u| + C = \frac{1}{3} \ln(3x^2 + 1) + C \]
\end{solution}

\question
\[ \int \sqrt{e^{\frac{3}{2}x} + e^x} \, dx \]
\begin{solution}
Factor $e^x$ out of the square root:
\[ \int \sqrt{e^x(e^{x/2} + 1)} \, dx = \int e^{x/2} \sqrt{e^{x/2} + 1} \, dx \]
Let $u = e^{x/2} + 1$, then $du = \frac{1}{2} e^{x/2} \, dx \implies 2 \, du = e^{x/2} \, dx$:
\[ \int 2\sqrt{u} \, du = \frac{4}{3} u^{3/2} + C = \frac{4}{3} \left(e^{x/2} + 1\right)^{3/2} + C \]
\end{solution}

\question
\[ \int_0^9 \sqrt{\frac{9}{x} - 1} \, dx \]
\begin{solution}
Use the substitution $x = 9 \sin^2(\theta)$, so $dx = 18 \sin(\theta) \cos(\theta) \, d\theta$.
The limits change as follows: $x = 0 \to \theta = 0$ and $x = 9 \to \theta = \frac{\pi}{2}$.
Substituting into the integrand yields:
\[ \sqrt{\frac{9}{9\sin^2(\theta)} - 1} = \sqrt{\csc^2(\theta) - 1} = \cot(\theta) \]
Now evaluate the transformed integral:
\[ \int_0^{\pi/2} \cot(\theta) \cdot 18 \sin(\theta) \cos(\theta) \, d\theta = 18 \int_0^{\pi/2} \cos^2(\theta) \, d\theta \]
Using the identity $\cos^2(\theta) = \frac{1 + \cos(2\theta)}{2}$:
\[ 9 \int_0^{\pi/2} (1 + \cos(2\theta)) \, d\theta = 9 \left[ \theta + \frac{\sin(2\theta)}{2} \right]_0^{\pi/2} = 9 \left( \frac{\pi}{2} \right) = \frac{9\pi}{2} \]
\end{solution}

\question
\[ \int_0^1 \frac{e^{\sqrt[3]{x}}}{\sqrt[3]{x}} \, dx \]
\begin{solution}
Let $u = \sqrt[3]{x} = x^{1/3}$, so $x = u^3$ and $dx = 3u^2 \, du$.
The limits remain $u = 0$ to $u = 1$:
\[ \int_0^1 \frac{e^u}{u} (3u^2 \, du) = 3 \int_0^1 u e^u \, du \]
Use integration by parts with $w = u, dv = e^u \, du \implies dw = du, v = e^u$:
\[ 3 \left[ u e^u - e^u \right]_0^1 = 3 \left( (e - e) - (0 - 1) \right) = 3 \]
\end{solution}

\question
\[ \int \left(e^x - \frac{e^x}{x}\right) \left(e^{\frac{1}{x}} + \frac{e^{\frac{1}{x}}}{x}\right) \, dx \]
\begin{solution}
Factor $e^x$ from the first factor and $e^{1/x}$ from the second factor:
\[ e^x \left(1 - \frac{1}{x}\right) e^{\frac{1}{x}} \left(1 + \frac{1}{x}\right) = e^{x + \frac{1}{x}} \left(1 - \frac{1}{x^2}\right) \]
Notice that if $u = x + \frac{1}{x}$, then $du = \left(1 - \frac{1}{x^2}\right) \, dx$.
The integral simplifies to:
\[ \int e^u \, du = e^u + C = e^{x + \frac{1}{x}} + C \]
\end{solution}

\question
\[ \int \frac{x+1}{\sqrt{x+4}\sqrt{x-2}} \, dx \]
\begin{solution}
Combine the radicals in the denominator:
\[ \sqrt{x+4}\sqrt{x-2} = \sqrt{(x+4)(x-2)} = \sqrt{x^2 + 2x - 8} \]
Notice that $\frac{d}{dx}(x^2 + 2x - 8) = 2x + 2 = 2(x + 1)$.
Let $u = x^2 + 2x - 8$, then $du = 2(x + 1) \, dx \implies (x+1) \, dx = \frac{1}{2} du$:
\[ \int \frac{\frac{1}{2}}{\sqrt{u}} \, du = \sqrt{u} + C = \sqrt{x^2 + 2x - 8} + C = \sqrt{(x+4)(x-2)} + C \]
\end{solution}

\question
\[ \int_{-1}^1 (x\sin(x) + x\cos(x))^2 \, dx \]
\begin{solution}
Expand the integrand:
\[ (x\sin(x) + x\cos(x))^2 = x^2 (\sin^2(x) + 2\sin(x)\cos(x) + \cos^2(x)) = x^2 (1 + \sin(2x)) = x^2 + x^2\sin(2x) \]
Since the integration domain $[-1, 1]$ is symmetric about $0$:
\begin{itemize}
    \item $x^2\sin(2x)$ is an odd function, so its integral over $[-1, 1]$ is $0$.
    \item $x^2$ is an even function.
\end{itemize}
\[ \int_{-1}^1 x^2 \, dx = 2 \int_0^1 x^2 \, dx = 2 \left[ \frac{x^3}{3} \right]_0^1 = \frac{2}{3} \]
\end{solution}

\question
\[ \int \frac{\sin(1-x) + \cos(1+x)}{\sin(1+x) + \cos(1-x)} \, dx \]
\begin{solution}
Let $u = \sin(1+x) + \cos(1-x)$.
Differentiate $u$ with respect to $x$:
\[ du = \left( \cos(1+x) - \sin(1-x)(-1) \right) \, dx = (\cos(1+x) + \sin(1-x)) \, dx \]
Notice that $du$ matches the numerator exactly. Thus:
\[ \int \frac{1}{u} \, du = \ln|u| + C = \ln|\sin(1+x) + \cos(1-x)| + C \]
\end{solution}

\question
\[ \int \frac{1}{\text{sech}(x) + \tanh(x)} \, dx \]
\begin{solution}
Rewrite in terms of $\sinh(x)$ and $\cosh(x)$:
\[ \text{sech}(x) + \tanh(x) = \frac{1}{\cosh(x)} + \frac{\sinh(x)}{\cosh(x)} = \frac{1 + \sinh(x)}{\cosh(x)} \]
The integral becomes:
\[ \int \frac{\cosh(x)}{1 + \sinh(x)} \, dx \]
Let $u = 1 + \sinh(x)$, then $du = \cosh(x) \, dx$:
\[ \int \frac{1}{u} \, du = \ln|u| + C = \ln(1 + \sinh(x)) + C \]
\end{solution}

\question
\[ \int_{\frac{\pi}{4}}^{\frac{\pi}{2}} \log\left( (\sin x)^{\tan x} (\sec x)^{\cot x} \right) \, dx \]
\begin{solution}
Using properties of logarithms and $\sec(x) = \frac{1}{\cos(x)}$:
\[ \log\left( (\sin x)^{\tan x} (\sec x)^{\cot x} \right) = \tan(x)\log(\sin x) - \cot(x)\log(\cos x) \]
Notice by product rule that:
\[ \frac{d}{dx} \left[ \log(\sin x)\log(\cos x) \right] = \cot(x)\log(\cos x) - \tan(x)\log(\sin x) \]
Therefore, the integrand is $-\frac{d}{dx} \left[ \log(\sin x)\log(\cos x) \right]$.
Evaluating the antiderivative at the limits:
\[ \left[ -\log(\sin x)\log(\cos x) \right]_{\pi/4}^{\pi/2} \]
As $x \to \frac{\pi}{2}^-$, $\sin(x) \to 1 \implies \log(\sin x) \to 0$, making the upper limit $0$.
At $x = \frac{\pi}{4}$: $\sin\left(\frac{\pi}{4}\right) = \cos\left(\frac{\pi}{4}\right) = \frac{1}{\sqrt{2}} = 2^{-1/2}$.
\[ 0 - \left( -\log\left(2^{-1/2}\right) \log\left(2^{-1/2}\right) \right) = \left( -\frac{1}{2}\log 2 \right)^2 = \frac{1}{4}(\log 2)^2 \]
\end{solution}

\question
\[ \lim_{n \to \infty} n \int_1^{n^2} \frac{\sin\left(\frac{x}{n}\right)}{e^x} \, dx \]
\begin{solution}
Let $t = \frac{1}{n}$. As $n \to \infty$, $t \to 0^+$. The limit expression becomes:
\[ \lim_{t \to 0^+} \frac{1}{t} \int_1^{1/t^2} e^{-x} \sin(tx) \, dx \]
Using the standard antiderivative $\int e^{-ax}\sin(bx) \, dx = \frac{e^{-ax}}{a^2+b^2}(-a\sin(bx) - b\cos(bx))$ with $a=1, b=t$:
\[ \int_1^{1/t^2} e^{-x} \sin(tx) \, dx = \left[ \frac{e^{-x}}{1+t^2}(-\sin(tx) - t\cos(tx)) \right]_1^{1/t^2} \]
As $t \to 0^+$, $e^{-1/t^2} \to 0$, so the upper bound vanishes. Evaluating at the lower bound gives:
\[ \frac{e^{-1}}{1+t^2}(\sin(t) + t\cos(t)) \]
Taking the limit:
\[ \lim_{t \to 0^+} \frac{e^{-1}}{1+t^2} \left( \frac{\sin(t)}{t} + \cos(t) \right) = e^{-1} \cdot 1 \cdot (1 + 1) = \frac{2}{e} \]
\end{solution}

\question
\[ \int_{-\frac{\pi}{2}}^{\frac{\pi}{2}} \frac{1 + \sin^{2023}(x)}{1 + \tan^{2022}(x)} \, dx \]
\begin{solution}
Split the integral into two parts:
\[ \int_{-\pi/2}^{\pi/2} \frac{\sin^{2023}(x)}{1 + \tan^{2022}(x)} \, dx + \int_{-\pi/2}^{\pi/2} \frac{1}{1 + \tan^{2022}(x)} \, dx \]
The integrand $\frac{\sin^{2023}(x)}{1 + \tan^{2022}(x)}$ is an odd function because $\sin^{2023}(-x) = -\sin^{2023}(x)$ and $\tan^{2022}(-x) = \tan^{2022}(x)$, so its integral over $[-\frac{\pi}{2}, \frac{\pi}{2}]$ is $0$.

For the remaining even part $I = \int_{-\pi/2}^{\pi/2} \frac{1}{1 + \tan^{2022}(x)} \, dx = 2 \int_0^{\pi/2} \frac{1}{1 + \tan^{2022}(x)} \, dx$:
Apply King's Property ($x \to \frac{\pi}{2} - x$):
\[ \int_0^{\pi/2} \frac{1}{1 + \tan^{2022}(x)} \, dx = \int_0^{\pi/2} \frac{1}{1 + \cot^{2022}(x)} \, dx = \int_0^{\pi/2} \frac{\tan^{2022}(x)}{1 + \tan^{2022}(x)} \, dx \]
Adding both integral forms:
\[ 2 \int_0^{\pi/2} \frac{1}{1 + \tan^{2022}(x)} \, dx = \int_0^{\pi/2} 1 \, dx = \frac{\pi}{2} \]
Thus, $I = \frac{\pi}{2}$.
\end{solution}

\question
\[ \int \frac{\log(x+1) - \log(x)}{x(x+1)} \, dx \]
\begin{solution}
Using partial fraction decomposition, $\frac{1}{x(x+1)} = \frac{1}{x} - \frac{1}{x+1}$.
Let $u = \log(x+1) - \log(x)$, then:
\[ du = \left( \frac{1}{x+1} - \frac{1}{x} \right) \, dx = -\frac{1}{x(x+1)} \, dx \]
The integral becomes:
\[ \int -u \, du = -\frac{u^2}{2} + C = -\frac{1}{2} (\log(x+1) - \log(x))^2 + C \]
\end{solution}

\question
\[ \int e^{x(e^x + 1)} (x + 1) \, dx \]
\begin{solution}
Rewrite the exponent as $x(e^x + 1) = x e^x + x$. The integral becomes:
\[ \int e^{x e^x + x} (x + 1) \, dx = \int e^{x e^x} e^x (x + 1) \, dx \]
Let $u = x e^x$, then $du = (e^x + x e^x) \, dx = e^x(x + 1) \, dx$:
\[ \int e^u \, du = e^u + C = e^{x e^x} + C \]
\end{solution}

\question
\[ \int (\log(x)\sin(x) + x\cos(x)\log(x)) \, dx \]
\begin{solution}
Factor $\log(x)$ out of the integrand:
\[ \int \log(x) (\sin(x) + x\cos(x)) \, dx \]
Observe that $\frac{d}{dx}(x\sin(x)) = \sin(x) + x\cos(x)$.
Apply integration by parts with $u = \log(x)$ and $dv = (\sin(x) + x\cos(x)) \, dx \implies v = x\sin(x), du = \frac{1}{x} \, dx$:
\[ x\log(x)\sin(x) - \int x\sin(x) \cdot \frac{1}{x} \, dx = x\log(x)\sin(x) - \int \sin(x) \, dx \]
\[ = x\log(x)\sin(x) + \cos(x) + C \]
\end{solution}

\question
\[ \int \frac{\log(x) + 2}{x\log(x) + x + 1} \, dx \]
\begin{solution}
Let $u = x\log(x) + x + 1$. Differentiating $u$ using the product rule:
\[ du = \left( \log(x) + x \cdot \frac{1}{x} + 1 \right) \, dx = (\log(x) + 2) \, dx \]
The numerator is precisely $du$. Substituting into the integral:
\[ \int \frac{1}{u} \, du = \ln|u| + C = \ln|x\log(x) + x + 1| + C \]
\end{solution}

\question
\[ \int \sqrt{\sec(x-7)} \tan(x-7) \, dx \]
\begin{solution}
Multiply the top and bottom by $\sqrt{\sec(x-7)}$:
\[ \int \frac{\sec(x-7)\tan(x-7)}{\sqrt{\sec(x-7)}} \, dx \]
Let $u = \sec(x-7)$, then $du = \sec(x-7)\tan(x-7) \, dx$:
\[ \int \frac{1}{\sqrt{u}} \, du = 2\sqrt{u} + C = 2\sqrt{\sec(x-7)} + C \]
\end{solution}

\question
\[ \int_{-\infty}^{\infty} (\cosh(x) - \sinh(x))^x \, dx \]
\begin{solution}
Recall that $\cosh(x) - \sinh(x) = e^{-x}$.
The integrand simplifies to:
\[ (\cosh(x) - \sinh(x))^x = (e^{-x})^x = e^{-x^2} \]
Evaluating the standard Gaussian integral:
\[ \int_{-\infty}^{\infty} e^{-x^2} \, dx = \sqrt{\pi} \]
\end{solution}

\question
\[ \int \left(\frac{1}{(x+3)^2} - \frac{1}{(x+2)^2}\right) \frac{dx}{\sqrt{(x+2)(x+3)}} \]
\begin{solution}
Notice that:
\[ \frac{1}{x+2} - \frac{1}{x+3} = \frac{1}{(x+2)(x+3)} \implies \sqrt{\frac{1}{x+2} - \frac{1}{x+3}} = \frac{1}{\sqrt{(x+2)(x+3)}} \]
Let $u = \frac{1}{x+2} - \frac{1}{x+3}$. Differentiating gives:
\[ du = \left( -\frac{1}{(x+2)^2} + \frac{1}{(x+3)^2} \right) \, dx = \left( \frac{1}{(x+3)^2} - \frac{1}{(x+2)^2} \right) \, dx \]
The integral simplifies directly to:
\[ \int \sqrt{u} \, du = \frac{2}{3} u^{3/2} + C = \frac{2}{3} \left( \frac{1}{x+2} - \frac{1}{x+3} \right)^{3/2} + C = \frac{2}{3\left((x+2)(x+3)\right)^{3/2}} + C \]
\end{solution}

\end{questions}